\documentclass{article}
\usepackage{amsmath,amssymb}
\usepackage{xurl}
\usepackage[hidelinks]{hyperref}
% Shared commands for notes in docs/paper-gaps/.

\DeclareUnicodeCharacter{2080}{\ifmmode{}_{0}\else\textsubscript{0}\fi}
\DeclareUnicodeCharacter{2081}{\ifmmode{}_{1}\else\textsubscript{1}\fi}
\DeclareUnicodeCharacter{2082}{\ifmmode{}_{2}\else\textsubscript{2}\fi}
\DeclareUnicodeCharacter{2099}{\ifmmode{}_{n}\else\textsubscript{n}\fi}

\newcommand{\Lean}{\textsc{Lean}}
\ExplSyntaxOn
\NewDocumentCommand{\leanid}{m}{
  \ifmmode
    \textsf{\detokenize{#1}}
  \else
    \group_begin:
    \urlstyle{sf}
    \tl_set:Nn \l_tmpa_tl {#1}
    \tl_replace_all:Nnn \l_tmpa_tl {\_} {_}
    \exp_args:NV \path \l_tmpa_tl
    \group_end:
  \fi
}
\ExplSyntaxOff
\providecommand{\tr}{\operatorname{tr}}
\newcommand{\tnleanIssueBase}{https://github.com/LionSR/TNLean/issues/}
\newcommand{\tnleanPrBase}{https://github.com/LionSR/TNLean/pull/}
\newcommand{\tnleanDocBase}{https://sirui-lu.com/TNLean/docs/}
\newcommand{\ghissue}[1]{\href{\tnleanIssueBase#1}{GitHub issue \##1}}
\newcommand{\ghpr}[1]{\href{\tnleanPrBase#1}{PR \##1}}
\newcommand{\doclink}[1]{\href{\tnleanDocBase#1.html}{\path{#1}}}

% Dirac notation and the identity operator, as in blueprint/src/macros/common.tex.
% \providecommand keeps note-local definitions authoritative.
\usepackage{dsfont}
\providecommand{\ket}[1]{|#1\rangle}
\providecommand{\bra}[1]{\langle #1|}
\providecommand{\braket}[2]{\langle #1 | #2 \rangle}
\providecommand{\ketbra}[2]{|#1\rangle\!\langle #2|}
\providecommand{\Id}{\mathds{1}}
\providecommand{\one}{\mathds{1}}
\providecommand{\C}{\mathbb{C}}
\providecommand{\MN}[1]{M_{#1}(\C)}
\providecommand{\MD}{M_D(\C)}

% Verdict marker: \gapnote{<kind>}{<status>}. Kind is the policy
% classification (clarification, local-correction, scope-restriction,
% unfaithful, false-source, open-gap); status is open, wip (elimination
% underway), resolved, or historical. Machine-read by the site index;
% prints a verdict line.
\newcommand{\gapnote}[2]{\par\noindent\textbf{Verdict:} #1 (#2).\par}

\title{The CZX PEPS: the Boundary Action on a Chain of Legs}
\author{}
\date{2026-09-27}
\begin{document}
\maketitle
\gapnote{scope-restriction}{open}

\section{What the source states}
The on-site $\mathbb Z_2$ symmetry of the CZX model is generated by
$U_{CZX}=U_XU_{CZ}$, with $U_X=X_1X_2X_3X_4$ and
$U_{CZ}=CZ_{12}CZ_{23}CZ_{34}CZ_{41}$ on the four qubits of a site, and the
plaquette ground state is invariant under it \cite{ChenLiuWen2011CZX}.%
\footnote{Source traceability:
    \path{References/1106.4752/source/dDSPTmodel.tex}, lines 283--312.}
On a region with boundary, the two qubits of a boundary half plaquette are
constrained to the span of $\ket{00}$ and $\ket{11}$ and form one effective
spin, and projecting the symmetry of the region onto the support of its
reduced state gives the effective boundary symmetry
$\prod_i\tilde X_i\prod_i\widetilde{CZ}_{i,i+1}$, a matrix product unitary of
bond dimension two.\footnote{Source traceability:
    \path{References/1106.4752/source/dDSPTmodel.tex}, lines 330--345 and
    377--385.}
The review states that the CZX PEPS of bond dimension four is invariant under
this on-site symmetry and that its virtual action is the CZX matrix product
unitary \cite{Cirac2021Matrix}.\footnote{Source traceability:
    \path{Papers/2011.12127/TN-Review-main.tex}, lines 2502--2516 and
    2582--2587.}

\section{What is formalized}
With the bond orientation of the note on the CZX tensor, each virtual leg
carries two qubits $(a,b)$, and the following hold.
\begin{enumerate}
    \item The CZX PEPS on a torus of width and height at least three is
          invariant under $U_{CZX}$ applied at every site.
    \item For one site tensor $A$, applying $U_{CZX}$ to the physical index
          equals applying $V=(X\otimes X)\,CZ$, that is
          $V\ket{ab}=(-1)^{ab}\ket{\bar a\bar b}$, to each of the four
          virtual legs.
    \item Let a closed chain of $N$ legs carry the pairs
          $(c_m,c_{m+1})$, indices modulo $N$, so that neighbouring legs
          carry the same qubit of their common plaquette, and let $\iota$ be
          the resulting isometry from $N$ effective spins into the legs. Then
          $V^{\otimes N}\iota=\iota\,U_N$ with
          $U_N\ket{c}=(-1)^{\sum_mc_mc_{m+1}}\ket{\bar c}$, the matrix product
          unitary with $T^{0,1}=\ket0\bra+$, $T^{1,0}=\ket1\bra-$, and
          $\iota^\dagger V^{\otimes N}\iota=U_N$.
\end{enumerate}

\section{The restriction}
Item~3 is stated on the chain of legs subject to the plaquette constraint,
not on the boundary of a region of the torus PEPS: that the contraction of
the tensors of a region, as a vector on its boundary legs, lies in the range
of $\iota$ is not formalized. In the source this is the statement that the
reduced state of the region is supported on the effective spins. The
restriction concerns only the identification of the region with the chain;
items~1 and~2 carry no restriction beyond the torus size recorded
elsewhere.\footnote{Results: \leanid{TNLean.PEPS.onSiteOperator_czxOnSite_mulVec_stateCoeff_czxPEPS},
    \leanid{TNLean.PEPS.czxOnSite_mul_czxSiteTensor},
    \leanid{TNLean.PEPS.onSiteOperator_czxLegOperator_mul_czxBoundaryEmbedding},
    \leanid{TNLean.PEPS.czxBoundaryEmbedding_conjTranspose_mul_onSiteOperator_mul}.}

\section{Elimination plan}
Contract the tensors of a rectangular region of the torus with free
boundary legs, show from the plaquette constraints inside the region that
neighbouring boundary legs carry equal qubits at their common plaquette, so
the region's vector lies in the range of $\iota$, and combine with item~2
applied at every site of the region.

\bibliographystyle{alpha}
\bibliography{references}
\end{document}
